\subsection{Primary decomposition in graded modules}

PROPOSITION 5. Let \(\Delta\) be a torsion-free commutative group, A a graded Noetherian ring of type \(\Delta\) , M a graded A-module of type M, N a graded submodule of M and \(N = \bigcap_{i \in I} Q_i\) a primary decomposition of N in M.

\begin{enumerate}
\def\labelenumi{(\roman{enumi})}
\item
  Let \(\mathbf{Q}_i'\) be the largest graded submodule of \(\mathbf{M}\) contained in \(\mathbf{Q}_i\) . Then the \(\mathbf{Q}_i'\) are primary and \(\mathbf{N} = \bigcap_{i \in I} \mathbf{Q}_i'\) .
\item
  If the primary decomposition \(\mathbf{N} = \bigcap_{i\in I}\mathbf{Q}_i\) is reduced, so is the primary decomposition \(\mathbf{N} = \bigcap_{i\in I}\mathbf{Q}_i'\) and for all \(i\in \mathbf{I}\) the prime ideals corresponding to \(\mathbf{Q}_i\) and \(\mathbf{Q}_i'\) are equal.
\item
  If \(\mathbf{Q}_i\) corresponds to a prime ideal \(\mathfrak{p}_i\) which is a minimal element of \(\operatorname{Ass}(\mathbf{M} / \mathbf{N})\) , \(\mathbf{Q}_i\) is a graded submodule of \(\mathbf{M}\) .
\end{enumerate}

We have seen (no. 2, Proposition 4) that the \(Q_{i}^{\prime}\) are primary with respect to M and \(N \subset Q_{i}^{\prime} \subset Q_{i}\) , which proves (i). Proposition 4 of no. 2 also shows that the prime ideal \(p_{i}^{\prime}\) corresponding to \(Q_{i}^{\prime}\) is the largest graded ideal contained in the prime ideal \(p_{i}\) corresponding to \(Q_{i}\) . If the decomposition \(N = \bigcap_{i \in I} Q_{i}\) is reduced, \(p_{i} \in \text{Ass}(M/N)\) for all \(i\) ( \(\S 2\) , no. 3, Proposition 4), hence \(p_{i}\) is a graded ideal (no. 1, Proposition 1) and therefore \(p_{i}^{\prime} = p_{i}\) ; then \(\text{Ass}(M/N) = \bigcup_{i \in I} \{p_{i}^{\prime}\}\) ( \(\S 2\) , no. 3, Proposition 4), which proves that the decomposition \(N = \bigcap_{i \in I} Q_{i}^{\prime}\) is reduced ( \(\S 2\) , no. 3, Proposition 4). Finally, if \(p_{i}\) is a minimal element of \(\text{Ass}(M/N)\) , then \(p_{i}^{\prime} = p_{i}\) since \(p_{i}\) is graded (no. 1, Proposition 1), whence \(Q_{i}^{\prime} = Q_{i}\) by virtue of \(\S 2\) , no. 3, Proposition 5.

% semantic-fragments: legacy-input-seam
\relax{}
